%% kapa: exam=1 round=35 year=2024 period=2 form=A subject=accounting subject_ko=회계학 seq=22 number=62
%% source: Q-Net 감정평가사 기출문제 (cst00309) artlSeq=5213229 file=raw/1008/5213229/2교시 문제지 원본.pdf
%% conversion: kordoc 4.24.51 -> tools/gen1.py (auto); review=proofread-sonnet; flags=tab
\begin{question}{62}
(주)감평은 20×1년 초 유형자산인 기계장치를 ￦50,000에 취득(내용연수 5년, 잔존가치 ￦0, 정액법 상각)하여 사용하고 있다. 20×2년 중 자산손상의 징후를 발견하고 손상차손을 인식하였으나 20×3년 말 손상이 회복되었다고 판단하였다. 동 기계장치의 순공정가치와 사용가치가 다음과 같을 때, 20×2년 말 인식할 손상차손(A)과 20×3년 말 인식할 손상차손환입액(B)은? (단, 동 기계장치는 원가모형을 적용한다.)
\par{\small\centering\begin{tabular}{|c|c|c|}\hline
구분 & 순공정가치 & 사용가치 \tabularnewline\hline
20×2년 말 & ￦15,000 & ￦18,000 \tabularnewline\hline
20×3년 말 & 21,000 & 17,000 \tabularnewline\hline
\end{tabular}\par}
\smallskip
\noindent\begin{tabular}{@{}l@{\hspace{0.4em}}p{0.13\linewidth}p{0.13\linewidth}@{\hspace{2em}}l@{\hspace{0.4em}}p{0.13\linewidth}p{0.13\linewidth}@{}}
 & \uline{A} & \uline{B} & & \uline{A} & \uline{B} \tabularnewline
① & ￦12,000 & ￦8,000 & ② & ￦12,000 & ￦9,000 \tabularnewline
③ & ￦15,000 & ￦8,000 & ④ & ￦15,000 & ￦9,000 \tabularnewline
⑤ & ￦15,000 & ￦12,000 & & & \tabularnewline
\end{tabular}
\end{question}
\answer{1}
